% !TEX program = xelatex
\documentclass[11pt]{article}
\usepackage[a4paper,margin=23mm]{geometry}
\usepackage{amsmath,amssymb,amsthm,mathtools,booktabs,array,longtable}
\usepackage{fontspec,microtype,xurl}
\usepackage[colorlinks=true,linkcolor=blue,urlcolor=blue]{hyperref}
\newcommand{\R}{\mathbb R}
\newcommand{\norm}[1]{\left\lVert#1\right\rVert}
\newcommand{\ip}[2]{\left\langle#1,#2\right\rangle}
\newcommand{\rank}{\operatorname{rank}}
\newcommand{\ran}{\operatorname{ran}}
\newcommand{\Span}{\operatorname{span}}
\newcommand{\diag}{\operatorname{diag}}
\newcommand{\negind}{\operatorname{n}_{-}}
\newcommand{\code}[1]{\ifmmode\text{\small\nolinkurl{#1}}\else{\small\nolinkurl{#1}}\fi}
\newcommand{\src}[2]{\href{https://github.com/tchen-research/NLAlib/blob/b7a96fd6ade2aad77a1a9a038ddf2791deea440e/#1}{#2}}
\newtheorem{theorem}{Theorem}[section]
\newtheorem{lemma}[theorem]{Lemma}
\newtheorem{proposition}[theorem]{Proposition}
\newtheorem{corollary}[theorem]{Corollary}
\theoremstyle{remark}\newtheorem{remark}[theorem]{Remark}
\numberwithin{equation}{section}
\setlength{\emergencystretch}{3em}
\setcounter{tocdepth}{2}
\title{The six remaining matrix toolkit entries\\Operator derivations for their eight scaffold lemmas}
\author{Mathematical handoff for NLAlib}
\date{13 October 2026}
\begin{document}
\maketitle

\noindent The refreshed public catalogue at commit
\code{b7a96fd6ade2aad77a1a9a038ddf2791deea440e} has 108 matrix-toolkit
entries: 102 marked proved and six marked scaffold. These are not six
unproved mathematical theories: most of their constituent Lean theorems
are already present. This note targets the eight declarations still
containing a deferred proof, and derives their stated consequences.
It contains mathematics, not a new Lean implementation.

\begin{center}\small
\begin{tabular}{@{}p{.30\linewidth}p{.61\linewidth}@{}}
\toprule Catalogue entry & Remaining mathematical bridge\\\midrule
Principal angles & One-column totalized tangent identity.\\
Cholesky recursion & Positive-pivot prefix, zero trailing columns, then rank.\\
Gram--Schmidt & Normalized rank-one projections identify the two recursions.\\
Givens rotations & A column-support invariant for the accumulated factor.\\
Hessenberg reduction & A block recursion; implicit-Q uniqueness by oriented intertwining.\\
Sylvester inertia/Sturm & Symmetric LU becomes LDL$^\top$; signs of pivot products count inertia.\\
\bottomrule
\end{tabular}
\end{center}

All matrices are real, transposes are Euclidean adjoints, and matrix norms
are Euclidean operator norms. Indices start at zero to match the Lean
statements. A span over no vectors is zero, an empty product is one, and
an empty matrix has rank zero. Whenever total division is used, we state
explicitly the convention $0^{-1}=0$. It is a library convention, not a
geometric assertion about an infinite tangent.

\paragraph{Source and reuse boundary.}
The \src{atlas/atlas.json}{pinned catalogue} and
\src{atlas/declarations.json}{declaration index} were checked against the
source files linked in each section. Named reuse results below have
\code{sorry: false} in that public index; this is a source-level check,
not a fresh kernel audit. The local working checkout predates these
modules, so public availability does not imply that a theorem can be
imported into that checkout without updating it. The public snapshot
uses Lean 4.33.1 and Mathlib commit
\code{0df444a360eaa60ab8c11dca51a86af692955474}.

\tableofcontents

\section{Principal angles: reduce the remaining bridge to one scalar}
\label{sec:angle}
\paragraph{Exact remaining target.}
\code{NLAlib.sq_tanThetaMax_replicateCol_eq_tanSqAngle}, in
\src{NLAlib/Krylov/Spectral/PowerMethod.lean}{PowerMethod.lean}.
The full principal-angle characterization is already recorded as proved;
there is no reason to reconstruct an SVD or a CS decomposition for this
one-column bridge.

Let $u,x\in\R^N$, $\ip uu=1$, and $\alpha=\ip ux$. Write $[v]$ for
the one-column matrix associated with a vector $v$. The definitions are
\[
 \operatorname{tanThetaMax}([u],[x])
 =\norm{(I-uu^\top)[x]\,[\alpha]^{-1}},\qquad
 \operatorname{tanSqAngle}(u,x)
 =\frac{\ip uu\ip xx-\ip ux^2}{\ip ux^2}.
\]
Both use total inversion at zero.

\begin{lemma}[The one-column operator norm]\label{lem:column}
For every vector $v$, $\norm{[v]}=\norm v_2$.
\end{lemma}
\begin{proof}
As a linear map from $\R$ to $\R^N$, the matrix $[v]$ sends $t$ to $tv$.
Thus $\norm{tv}_2=|t|\norm v_2$, giving the upper bound. Evaluate at
$t=1$ for the reverse bound. This includes $v=0$ and $N=0$.
\end{proof}

\begin{theorem}[The exact totalized identity]
For $\ip uu=1$ and every $x$,
\[
 \operatorname{tanThetaMax}([u],[x])^2
     =\operatorname{tanSqAngle}(u,x).
\]
\end{theorem}
\begin{proof}
Put $r=x-\alpha u$. Finite bilinear expansion gives
\[
 \ip ur=0,\qquad \norm r_2^2=\norm x_2^2-\alpha^2.
\]
The overlap matrix is the $1\times1$ matrix $[\alpha]$ and its total
matrix inverse is $[\alpha^{-1}]$, also at $\alpha=0$. The residual
product is therefore $[\alpha^{-1}r]$. By Lemma~\ref{lem:column},
\[
 \operatorname{tanThetaMax}([u],[x])^2
 =|\alpha^{-1}|^2\norm r_2^2
 =\frac{\norm x_2^2-\alpha^2}{\alpha^2}
 =\operatorname{tanSqAngle}(u,x).
\]
For $\alpha=0$ both sides equal zero by total inversion; this branch
uses no cancellation of $\alpha$. For $\alpha\ne0$ it is ordinary division.
\end{proof}

\paragraph{Interpretation and edge cases.}
For $x\ne0$ and $\alpha\ne0$ this is the usual squared tangent of the
acute angle between the two lines. When $\alpha=0$ and $x\ne0$, the
geometric angle is $\pi/2$; the zero value is only the chosen totalization.
When $x=0$, both expressions still equal zero but there is no nonzero
second line. If $N=0$, the hypothesis $\ip uu=1$ is impossible, so the
theorem is vacuous rather than requiring an invented unit vector.

\paragraph{Reuse.}
The exact definitions are in
\src{NLAlib/Matrix/Subspaces/PrincipalAngles.lean}{PrincipalAngles.lean}
and \src{NLAlib/Krylov/Spectral/Ritz.lean}{Ritz.lean}. The general angle
package is \code{NLAlib.principal_angle_characterizations}.
For the remaining bridge prove or reuse the one-column operator norm,
the $1\times1$ inverse, and the displayed residual identity. Be careful:
Mathlib's unqualified \code{Matrix.norm_replicateCol} is a theorem about
its ambient matrix norm instance and is not automatically the Euclidean
operator-norm statement. Working with the map $t\mapsto tv$ avoids that
ambiguity completely.

\section{Pivoted Cholesky: establish the zero pattern before counting rank}
\label{sec:chol}
\paragraph{Exact remaining targets.}
\code{NLAlib.choleskyPivoted_apply_self_pos_iff_lt_rank} and
\code{NLAlib.choleskyPivoted_apply_eq_zero_of_rank_le}, in
\src{NLAlib/Matrix/Factorizations/Cholesky.lean}{Cholesky.lean}.
The factorization identity and positive-semidefinite correctness of the
recursion are already proved. The missing assertion is the precise number
and location of its nonzero columns.

\subsection{The scalar PSD observation behind stopping}
\begin{lemma}[A zero diagonal entry has a zero row and column]\label{lem:psdzero}
If $A\succeq0$ and $a_{ii}=0$, then $a_{ij}=a_{ji}=0$ for every $j$.
Consequently a PSD matrix with all diagonal entries zero is zero.
\end{lemma}
\begin{proof}
For all $t\in\R$,
\[
 0\le (e_i+t e_j)^\top A(e_i+t e_j)=2t a_{ij}+t^2a_{jj}.
\]
If $a_{ij}\ne0$, choose a sufficiently small nonzero $t$ with opposite
sign to $a_{ij}$; the linear term makes this expression negative.
Thus $a_{ij}=0$. Symmetry supplies the row/column equality.
\end{proof}
In particular, the largest diagonal of a nonzero PSD matrix is strictly
positive. Since all its diagonal entries are nonnegative, maximizing
their absolute values is the same pivot rule.

\subsection{One block step and a stronger induction invariant}
After the first symmetric swap, write
\[
 B=\begin{pmatrix}a&b^\top\\b&C\end{pmatrix},\qquad
 S=C-a^{-1}bb^\top.
\]
If $B\ne0$, the pivot rule makes $a>0$. Completing the square gives
\begin{equation}\label{eq:squares}
 \begin{pmatrix}t\\y\end{pmatrix}^{\!\top}
 B\begin{pmatrix}t\\y\end{pmatrix}
 =a(t+a^{-1}b^\top y)^2+y^\top S y.
\end{equation}
Hence $S\succeq0$. If $a=0$, maximality of the pivot forces $B=0$,
and the recursion continues on zero matrices with zero columns.

Suppose the recursive factorization of $S$ uses a trailing permutation
$P$ and lower triangular factor $L_S$, so $PSP^\top=L_SL_S^\top$.
The assembled factor is
\begin{equation}\label{eq:pivotL}
 L=\begin{pmatrix}\sqrt a&0\\Pb/\sqrt a&L_S\end{pmatrix}.
\end{equation}
This formula explains why subsequent pivots permute the lower part of
the already computed first column. It preserves the factorization:
\[
 \begin{pmatrix}1&0\\0&P\end{pmatrix}B
 \begin{pmatrix}1&0\\0&P^\top\end{pmatrix}=LL^\top.
\]
The equality follows by multiplying the four blocks. The zero case is
valid as well because $b=C=0$ and the total quotient is zero.

\begin{lemma}[Positive prefix and zero tail, without mentioning rank]
\label{lem:prefix}
For the pivoted Cholesky output of every PSD $N\times N$ matrix, there
is an integer $r\in\{0,\ldots,N\}$ such that
\[
 L_{jj}>0\ (j<r),\qquad L_{ij}=0\ (j\ge r,\ 0\le i<N).
\]
The factor is lower triangular and its diagonal is nonnegative.
\end{lemma}
\begin{proof}
Induct on $N$. At $N=0$ take $r=0$. At positive dimension, a zero maximal
pivot makes the whole matrix zero by Lemma~\ref{lem:psdzero}; induction
and the block formula give $L=0$, so take $r=0$.
Otherwise $a>0$. The PSD Schur complement has a recursive prefix length
$s$. Formula~\eqref{eq:pivotL} has positive first diagonal entry and
the positive diagonal entries of $L_S$ next, while every column after
column $s$ is zero (its top entry is zero and its trailing column is
zero). Take $r=1+s$. Lower triangularity and nonnegative diagonal also
follow from the same formula. No rank identity is used in the induction.
\end{proof}

\subsection{A finite rank certificate identifies the prefix length}
\begin{lemma}[Rank of the factor and of its Gram matrix]\label{lem:rankL}
A lower triangular matrix satisfying Lemma~\ref{lem:prefix} has rank $r$.
Every real matrix $L$ satisfies $\rank(LL^\top)=\rank L$.
\end{lemma}
\begin{proof}
Only its first $r$ columns can be nonzero, so the rank is at most $r$.
The first $r\times r$ principal block is triangular with determinant
$\prod_{j<r}L_{jj}>0$, so those columns are independent and the rank
is at least $r$. For $r=0$ the matrix is zero and the conclusion is direct.
Finally
\[
 LL^\top z=0\quad\Longleftrightarrow\quad L^\top z=0,
\]
because $z^\top LL^\top z=\norm{L^\top z}_2^2$. Rank--nullity and
$\rank L^\top=\rank L$ give the Gram-rank identity.
\end{proof}

\begin{theorem}[Both remaining Cholesky conclusions]
If $A\succeq0$ and the pivoted recursion returns $(\sigma,L)$, then
\[
 L_{ii}>0\quad\Longleftrightarrow\quad i<\rank A,
 \qquad
 j\ge\rank A\quad\Longrightarrow\quad L_{ij}=0\quad\text{for all }i.
\]
Thus it takes exactly $\rank A$ positive-pivot steps, and the first
$\rank A$ columns alone give the permuted Gram factorization of $A$.
\end{theorem}
\begin{proof}
The existing recursion theorem gives $P_\sigma A P_\sigma^\top=LL^\top$.
Permutations preserve rank. By Lemma~\ref{lem:rankL},
$\rank A=\rank(LL^\top)=\rank L=r$. Substitute this into
Lemma~\ref{lem:prefix}.
\end{proof}

\paragraph{Reuse and edge cases.}
Reuse \code{NLAlib.isPivotedCholesky_choleskyPivoted},
\code{NLAlib.abs_le_abs_absArgmax},
\code{NLAlib.posSemidef_pivotSchur}, and
\code{NLAlib.apply_eq_zero_of_posSemidef_of_apply_self_eq_zero}.
Mathlib has \code{Matrix.rank_self_mul_transpose} in
\code{Mathlib/LinearAlgebra/Matrix/Rank.lean}.
The useful new induction lemma is \emph{the prefix statement before rank}.
No rank-of-Schur-complement theorem is necessary. Ties in the maximum
diagonal can be broken arbitrarily. The zero matrix has $r=0$; a positive
definite matrix has $r=N$; dimension zero has no pivot and no exceptional
index to handle.

\section{Gram--Schmidt: identify the recursions through one rank-one projector}
\label{sec:gs}
\paragraph{Exact remaining target.}
\code{NLAlib.gramSchmidtCol_eq_gramSchmidtNormed}, in
\src{NLAlib/Matrix/Factorizations/GramSchmidt.lean}{GramSchmidt.lean}.
This target has \emph{no linear-independence hypothesis}. Both definitions
continue after a zero residual by assigning the zero vector.

\subsection{A normalization identity valid also at zero}
For $v$ in a real finite-dimensional inner-product space define
$\mathcal N(v)=\norm v^{-1}v$, with $\mathcal N(0)=0$.
\begin{lemma}\label{lem:rankone}
For every $v,y$,
\begin{equation}\label{eq:rankone}
 \ip{\mathcal N(v)}y\mathcal N(v)
 =\frac{\ip vy}{\norm v^2}v
 =P_{\R v}y.
\end{equation}
\end{lemma}
\begin{proof}
If $v=0$, every expression is zero, including projection onto the zero
subspace. Otherwise put $q=v/\norm v$. Then $\norm q=1$;
$\ip qyq$ belongs to $\R v$, and $y-\ip qyq$ is orthogonal to $v$.
These two properties uniquely characterize the orthogonal projection.
Bilinearity gives the scalar quotient in \eqref{eq:rankone}.
\end{proof}

Let $a_j$ be the columns of $A$. The normalized recursion in NLAlib is
\begin{equation}\label{eq:cgs}
 v_j=a_j-\sum_{i<j}\ip{q_i}{a_j}q_i,\qquad q_j=\mathcal N(v_j).
\end{equation}
Mathlib first defines the unnormalized recursion
\begin{equation}\label{eq:mathgs}
 g_j=a_j-\sum_{i<j}P_{\R g_i}a_j,
 \qquad \widehat q_j=\mathcal N(g_j).
\end{equation}

\begin{theorem}[Equality of the two recursions, including breakdown]
For every matrix $A$ and every column $j$, $v_j=g_j$ and $q_j=\widehat q_j$.
\end{theorem}
\begin{proof}
Use strong induction on $j$. Assuming $q_i=\mathcal N(g_i)$ for every
$i<j$, Lemma~\ref{lem:rankone} identifies each summand in
\eqref{eq:cgs} with its counterpart in \eqref{eq:mathgs}. Hence $v_j=g_j$.
Apply the same normalization map to obtain $q_j=\widehat q_j$.
The argument includes an empty predecessor sum and zero earlier vectors.
\end{proof}

\paragraph{The coordinate bridge.}
The map taking a coordinate vector $v:m\to\R$ to
\code{WithLp.toLp 2 v} preserves addition, finite sums, scalar
multiplication and inner products, and its norm is
$\sqrt{\sum_i v_i^2}$. Thus the theorem gives exactly
\[
 \code{gramSchmidtCol A j}
 =\bigl(\code{InnerProductSpace.gramSchmidtNormed}
       \ \R\ (a_i)_i\ j\bigr)\text{ in coordinates}.
\]
There is no need to prove QR uniqueness to identify the two definitions.

\subsection{Why the classical and modified algorithms agree}
This is already proved in the source, but the operator identity explains
the reuse. The nonzero earlier $q_i$ are orthonormal. Put $P_i=q_iq_i^\top$;
when $q_i=0$ this is the zero operator. For $i\ne j$, $P_iP_j=0$, so
finite induction gives
\begin{equation}\label{eq:projectorproduct}
 (I-P_{j-1})\cdots(I-P_0)=I-\sum_{i<j}P_i.
\end{equation}
The left side is the modified Gram--Schmidt residual operator and the
right side the classical one. Moreover
$\ip{q_i}{(I-P_{i-1})\cdots(I-P_0)a_j}=\ip{q_i}{a_j}$, so their stored
coefficients agree as well. This proves equality in exact arithmetic,
including the convention at breakdown.

The sum $\sum_{i<j}P_i$ projects onto the span of the earlier columns:
inductively, $a_i=v_i+\sum_{h<i}\ip{q_h}{a_i}q_h$ and $q_i$ is a scalar
multiple of $v_i$, so the two successive column spans agree. Consequently
$v_j=0$ exactly when $a_j$ is in the earlier column span. If no residual
vanishes, the $q_j$ form the graded orthonormal basis and the positive
diagonal entries of $R$ are $\norm{v_j}$.

\paragraph{Reuse.}
Mathlib's
\code{InnerProductSpace.gramSchmidt_def},
\code{InnerProductSpace.gramSchmidtNormed}, and
\code{Submodule.starProjection_singleton} expose precisely
\eqref{eq:mathgs} and \eqref{eq:rankone}.
The NLAlib results \code{gramSchmidtCol_eq},
\code{gramSchmidtResidual_eq_zero_iff}, and
\code{modifiedGramSchmidt_eq} already cover the other interfaces.
Empty column sets require no induction step. With an empty row set all
residuals and normalized vectors are zero; the bridge still holds.

\section{Givens QR: prove a column-support invariant}
\label{sec:givens}
\paragraph{Exact remaining target.}
\code{NLAlib.givensQRQ_invariant}, in
\src{NLAlib/Matrix/Factorizations/Givens.lean}{Givens.lean}.
Its three dependent structure-preservation results are already written
against this scaffold. The invariant holds for every input matrix,
not just Hessenberg inputs.

Write $F_j=\Span(e_0,\ldots,e_j)$. A matrix is upper Hessenberg exactly
when its column $j$ belongs to $F_{j+1}$ (truncated at the full ambient
space). Let $G_k$ be any orthogonal rotation acting only in the plane
$\Span(e_k,e_{k+1})$, and set $G_k=I$ when $k+1\ge N$. Define
\[
 Q_0=I,\qquad Q_{k+1}=Q_kG_k^\top.
\]
The rotations may depend on all previous steps; the support proof does
not depend on their angles or on how they were chosen.

\begin{theorem}[The exact accumulated-factor invariant]\label{thm:givens}
For every $k\ge0$, $Q_k$ is upper Hessenberg and
$Q_ke_j=e_j$ for every coordinate index $j>k$.
\end{theorem}
\begin{proof}
At $k=0$ the statement is immediate. If $k+1\ge N$, the next multiplier
is the identity. Upper Hessenberg shape persists, and the new condition
on $j>k+1$ follows from the old one.

Otherwise right multiplication changes only columns $k,k+1$. Write the
two columns of the $2\times2$ block of $G_k^\top$ as $(c,s)$ and
$(-s,c)$, so
\begin{align*}
 Q_{k+1}e_k&=cQ_ke_k+sQ_ke_{k+1},\\
 Q_{k+1}e_{k+1}&=-sQ_ke_k+cQ_ke_{k+1}.
\end{align*}
The old Hessenberg condition gives $Q_ke_k\in F_{k+1}$, and the old
identity-column condition gives $Q_ke_{k+1}=e_{k+1}\in F_{k+1}$.
Both new columns therefore belong to $F_{k+1}$. This is the exact
Hessenberg requirement for column $k$ and a stronger requirement than
needed for column $k+1$. Every other column is unchanged. In particular,
all columns $j>k+1$ are still identity columns. This closes the induction.
\end{proof}

\paragraph{Equivalent entrywise induction.}
If $a>k+1$, the two entries mixed in row $a$ are zero:
$(Q_k)_{ak}=0$ by Hessenberg shape and
$(Q_k)_{a,k+1}=\delta_{a,k+1}=0$ by the identity-column condition.
The transformed entries are therefore zero. All untouched columns use
the old invariant. This proves the exact pair of clauses in the Lean
statement using only finite sums with two surviving terms; it avoids
transporting a row lemma through transposes.

\begin{corollary}[The QR step preserves shape even at singular matrices]
Let $H$ be upper Hessenberg and let $H=QR$ be the output of the ordered
Givens recursion. Then $RQ$ is upper Hessenberg. If $H$ is symmetric,
$RQ$ is symmetric tridiagonal. No invertibility hypothesis on $R$ is needed.
\end{corollary}
\begin{proof}
The recursion already gives upper triangular $R$ and orthogonal $Q$;
Theorem~\ref{thm:givens} gives upper Hessenberg $Q$. For $i>j+1$,
\[
 (RQ)_{ij}=\sum_\ell R_{i\ell}Q_{\ell j}=0:
\]
terms with $\ell<i$ have $R_{i\ell}=0$, and those with $\ell\ge i>j+1$
have $Q_{\ell j}=0$. Also $RQ=Q^\top H Q$, so it is symmetric when
$H$ is. A symmetric upper Hessenberg matrix is tridiagonal.
\end{proof}

\paragraph{Reuse and small dimensions.}
Use \code{NLAlib.isQR_givensQR},
\code{NLAlib.isUpperTriangular_givensQR_snd},
\code{NLAlib.IsUpperHessenberg.isUpperTriangular_mul}, and
\code{NLAlib.isSymm_mul_of_isQR}. For $N=0,1$ the rotation list is empty
and $Q=I$. A zero two-entry pivot gives an identity rotation, so no
division-by-zero assumption enters the support argument.

\section{Hessenberg reduction: block recursion and oriented intertwining}
\label{sec:hess}
\paragraph{Exact remaining targets.}
\code{NLAlib.isUpperHessenberg_hessenbergH} and
\code{NLAlib.IsHessenbergForm.exists_sign}, in
\src{NLAlib/Matrix/Factorizations/Hessenberg.lean}{Hessenberg.lean}.
Orthogonality of the constructed factor, its fixed first column, and the
number of reflectors are already proved. Preserve this actual algorithm;
an unrelated existential Schur factorization would not prove its scaffold.

\subsection{One padding identity handles the recursive algorithm}
For $v\in\R^N$, the total Householder map is
\[
 H_v=I-2\frac{vv^\top}{v^\top v},\qquad H_0=I.
\]
It is symmetric and orthogonal. If $v=(0,w)$, direct block multiplication
gives
\begin{equation}\label{eq:padding}
 H_{(0,w)}=\diag(1,H_w).
\end{equation}
Consequently products of padded reflectors equal the padded product;
this follows by induction from
$\diag(1,U)\diag(1,V)=\diag(1,UV)$, including the empty product.

For completeness, the usual reflector taking $x$ to a multiple of $e_0$
works with zero tails as well. Put $r=\norm x_2$,
$s=1$ if $x_0\ge0$ and $s=-1$ otherwise, and $w=x+sr e_0$.
If $x=0$, $w=0$ and the reflector is the identity. Otherwise
\[
 w^\top x=r(r+sx_0),\qquad w^\top w=2r(r+sx_0)>0,
\]
so $H_wx=x-w=-sr e_0$. This is the scalar algebra used by the
existing \code{householderVector} interface.

\begin{theorem}[Correctness of the concrete Hessenberg recursion]
Let $Q(A)$ be the product of the reflectors in the existing recursion.
Then $Q(A)^\top A Q(A)$ is upper Hessenberg for every square $A$.
The construction uses $\max(N-2,0)$ reflectors in dimension $N$.
\end{theorem}
\begin{proof}
Use induction on $N$. Dimensions $0,1,2$ already have Hessenberg shape
and the product is the identity. For $N\ge3$, write
\[
 A=\begin{pmatrix}a&z^\top\\x&C\end{pmatrix},\quad
 P=\diag(1,H_w),\quad H_wx=\beta e_0,\quad B=PAP.
\]
Since $Pe_0=e_0$, the first column of $B$ is
$(a,\beta,0,\ldots,0)^\top$. Let $C'$ be the trailing $(N-1)\times(N-1)$
block of $B$, and put $U=Q(C')$. The recursive reflector-list definition
and \eqref{eq:padding} give the exact factor identity
\begin{equation}\label{eq:hessrec}
 Q(A)=P\diag(1,U).
\end{equation}
Every recursive reflector fixes its first coordinate, so $Ue_0=e_0$.
Orthogonality then also gives $U^\top e_0=e_0$. Consequently
\[
 Q(A)^\top A Q(A)
 =\diag(1,U^\top)B\diag(1,U)
 =\begin{pmatrix}a&*\\\beta e_0&U^\top C'U\end{pmatrix}.
\]
The trailing block is Hessenberg by induction and the first column has
the required support. Those two assertions are exactly upper Hessenberg
shape for the whole matrix. Counting the first reflector plus those in
dimension $N-1$ gives $1+(N-3)=N-2$.
\end{proof}

\paragraph{Immediate consequences and reuse.}
The displayed factor is orthogonal, so $A=QHQ^\top$. If $A=A^\top$,
then $H=H^\top$, and upper Hessenberg shape makes it tridiagonal.
Use \code{NLAlib.householder_vecCons_zero},
\code{NLAlib.householder_mulVec_householderVector},
\code{NLAlib.hessenbergQ_mulVec_single_zero}, and
\code{NLAlib.hessenbergQ_transpose_mul_self}.
The only list lemma required beyond these is the finite product identity
behind \eqref{eq:hessrec}. No continuity, spectral theorem, or nonzero-tail
assumption is involved.

\subsection{The implicit-Q theorem: put the known unreduced matrix on the right}
\begin{lemma}[Triangularity from an oriented intertwining equation]
\label{lem:intertwine}
Suppose $H,G\in\R^{N\times N}$ are upper Hessenberg, with
$h_{k+1,k}\ne0$ for $0\le k<N-1$. If
\[
 GW=WH,\qquad We_0=e_0,
\]
then $W$ is upper triangular. Here $N\ge1$.
\end{lemma}
\begin{proof}
Write $w_j=We_j$. Inductively prove $w_j\in F_j$.
The base is $w_0=e_0$. Comparing column $k$ in the intertwining equation
gives
\begin{equation}\label{eq:intertwine}
 h_{k+1,k}w_{k+1}=Gw_k-\sum_{i=0}^k h_{ik}w_i.
\end{equation}
Because $G$ is upper Hessenberg, $G F_k\subseteq F_{k+1}$. Every term
on the right therefore lies in $F_{k+1}$. Divide only by the known
nonzero number $h_{k+1,k}$. It follows that $w_{k+1}\in F_{k+1}$.
\end{proof}

\begin{lemma}[Orthogonal and upper triangular means diagonal signs]
\label{lem:orthtri}
If $W^\top W=I$ and $W$ is square upper triangular, then
$W=\diag(d)$ with $d_j\in\{1,-1\}$.
\end{lemma}
\begin{proof}
Column zero is $d_0e_0$ and its norm is one. Orthogonality with the later
columns makes their zeroth entries zero. The trailing block is again
upper triangular and orthogonal. Induct on dimension. At dimension zero
there is just the empty diagonal matrix.
\end{proof}

\begin{theorem}[Implicit Q under exactly the scaffold hypotheses]
Let $N\ge1$ and
\[
 A=Q_1H_1Q_1^\top=Q_2H_2Q_2^\top,
\]
where both $Q_i$ are orthogonal and both $H_i$ upper Hessenberg. Assume
only $H_1$ is unreduced, and $Q_1e_0=Q_2e_0$. Then there is a diagonal
sign matrix $D$, with $D_{00}=1$, such that
\[
 Q_2=Q_1D,\qquad H_2=DH_1D.
\]
\end{theorem}
\begin{proof}
Choose $W=Q_2^\top Q_1$. Then $H_2W=WH_1$, $We_0=e_0$, and $W$ is
orthogonal. Lemmas~\ref{lem:intertwine} and \ref{lem:orthtri} give
$W=D$. Since $D^2=I$ and its first column is $e_0$, $D_{00}=1$.
Now $Q_1=Q_2D$ gives $Q_2=Q_1D$, and conjugating $A$ gives the asserted
formula for $H_2$.
\end{proof}

\paragraph{Why this orientation matters.}
If one instead writes $W=Q_1^\top Q_2$ and compares columns of
$H_1W=WH_2$, the multiplier of $w_{k+1}$ is $(H_2)_{k+1,k}$, which
was not assumed nonzero. The source's informal sketch mixes these two
orientations. Choosing $W=Q_2^\top Q_1$ fixes the argument without
strengthening the theorem. In fact the conclusion implies that $H_2$
is unreduced as well.

\subsection{The Arnoldi interpretation stays a consequence}
Write $q_j=Q_1e_j$ and
$\mathcal K_{j+1}=\Span(q_0,Aq_0,\ldots,A^jq_0)$. Since $AQ_1=Q_1H_1$,
\[
 Aq_j=\sum_{i=0}^{j+1}(H_1)_{ij}q_i.
\]
Induction, solving for $q_{j+1}$ using the nonzero subdiagonal, gives
$\mathcal K_{j+1}=\Span(q_0,\ldots,q_j)$. Thus these columns are the
graded orthonormal Krylov basis, up to the chosen signs. A convention
of positive subdiagonal entries removes those signs once $q_0$ is fixed.
The already-proved \code{NLAlib.ArnoldiDecomp.exists_sign} expresses
related uniqueness, but importing the Krylov layer into the matrix layer
would violate the dependency direction. Lemma~\ref{lem:intertwine}
is the short matrix-only proof needed here. If a subdiagonal vanishes,
the induction stops at that index; unrestricted uniqueness of the
remaining columns is not claimed.

\section{Sturm counting: symmetric LU and signs of finite products}
\label{sec:sturm}
\paragraph{Exact remaining target.}
\code{NLAlib.card_eigenvalues_lt_eq_card_sign_changes}, in
\src{NLAlib/Matrix/Spectral/Inertia.lean}{Inertia.lean}.
This is stated for \emph{every real symmetric matrix} whose shifted
leading principal minors are nonzero, not only for tridiagonal matrices.
The law of inertia and the eigenvalue count from an LDL$^\top$
factorization are already proved.

Let $A=A^\top\in\R^{N\times N}$, $x\in\R$, and $M=A-xI$. Define
\[
 p_0=1,\qquad p_k=\det M_{0:k-1,\,0:k-1}\quad(1\le k\le N).
\]
The exact scaffold assumes $p_k\ne0$ for $1\le k\le N$.

\subsection{Extract LDL from the already available LU}
\begin{lemma}[Symmetric LU needs no second elimination algorithm]
\label{lem:luld}
If $M=M^\top=LU$, with $L$ unit lower triangular and $U$ upper triangular,
then $M=LDL^\top$ with $D=\diag(U_{00},\ldots,U_{N-1,N-1})$.
\end{lemma}
\begin{proof}
The unit lower triangular matrix $L$ is invertible. Set
\[
 B=L^{-1}ML^{-\top}=UL^{-\top}.
\]
The left expression is symmetric and the right expression is upper
triangular. Hence $B$ is diagonal. In the diagonal entry of the product
$UL^{-\top}$, upper triangularity forces the intermediate index to equal
the diagonal index, and $L^{-\top}$ has diagonal one. Therefore
$B_{ii}=U_{ii}$. Multiplying back gives the factorization.
\end{proof}

\begin{theorem}[The nonzero-leading-minor sign count]\label{thm:sturm}
If $p_k\ne0$ for all $1\le k\le N$, then, with eigenvalues counted
with multiplicity,
\begin{equation}\label{eq:sturm}
 \#\{i:\lambda_i(A)<x\}
    =\#\{k\in\{0,\ldots,N-1\}:p_kp_{k+1}<0\}.
\end{equation}
\end{theorem}
\begin{proof}
Nonzero leading minors give successful ordinary LU elimination and
$M=LU$. Lemma~\ref{lem:luld} gives $M=LDL^\top$ with $d_k=U_{kk}$.
The already-established law of inertia gives
\[
 \#\{i:\lambda_i(A)<x\}=\negind(M)=\#\{k:d_k<0\}.
\]
The leading block of $LU$ is the product of its leading blocks, since
the triangular factors eliminate all terms beyond the block. Taking
determinants yields the finite products
\[
 p_k=\prod_{i<k}d_i,\qquad p_{k+1}=p_kd_k.
\]
Thus
\begin{equation}\label{eq:pivotsign}
 p_kp_{k+1}=p_k^2d_k,
 \qquad p_k^2>0,
 \qquad d_k<0\ \Longleftrightarrow\ p_kp_{k+1}<0.
\end{equation}
Counting identical finite predicates proves \eqref{eq:sturm}. This uses
no rational inequality: division by $p_k$ can be omitted entirely.
For $N=0$ both counts are zero and $p_0=1$.
\end{proof}

\paragraph{Minimal reusable ingredients.}
Use \code{NLAlib.luSucceeds_iff_forall_det_ne_zero} and
\code{NLAlib.isLU_lu} in
\src{NLAlib/Matrix/Factorizations/LU.lean}{LU.lean}, followed by
\code{NLAlib.IsLU.det_leadingSubmatrix} and
\code{NLAlib.card_eigenvalues_lt_eq_card_neg_of_eq_mul_diagonal_mul}.
The one definition bridge says that \code{leadingMinor M k} equals
the determinant of \code{leadingSubmatrix M k hk} when $k\le N$.
The only additional matrix lemma is Lemma~\ref{lem:luld}.
Repeated eigenvalues cause no problem: the inertia theorem already
counts multiplicities. Under these hypotheses $p_N\ne0$, so $x$ is
not an eigenvalue of $A$.

\subsection{Tridiagonal recurrence is a finite determinant identity}
For a symmetric tridiagonal matrix with diagonal $a_0,\ldots,a_{N-1}$
and off-diagonal $b_0,\ldots,b_{N-2}$, expansion along the last row gives
\begin{equation}\label{eq:three}
 p_0(x)=1,\quad p_1(x)=a_0-x,\quad
 p_{k+1}(x)=(a_k-x)p_k(x)-b_{k-1}^2p_{k-1}(x).
\end{equation}
Indeed the diagonal term is $(a_k-x)p_k$; the only other cofactor
contains another factor $b_{k-1}$ and determinant $p_{k-1}$, with the
negative cofactor sign. Thus evaluating the list by the three-term
recursion and counting sign changes gives the spectral count whenever
the terms are nonzero. No root separation or eigenvalue continuity
theorem is needed.

\subsection{A stronger zero-safe Sturm rule by two-by-two elimination}
This optional extension supplies the boundary cases used in bisection.
It is not needed to prove the current nonzero-minor scaffold.
For a finite real list $z$, let $V(z)$ be its number of sign changes
\emph{after deleting all zero terms}.

\begin{theorem}[Unreduced tridiagonal count at every shift]\label{thm:zerosturm}
If $T=T^\top$ is tridiagonal and every $b_j\ne0$, then for every real $x$,
\[
 \#\{i:\lambda_i(T)<x\}=V(p_0(x),p_1(x),\ldots,p_N(x)).
\]
This includes $x$ equal to an eigenvalue and any vanishing interior minor.
\end{theorem}
\begin{proof}
We prove by dimension induction that for any such tridiagonal matrix
$M$, $\negind(M)$ equals the variation of its leading determinants.
Dimension zero gives the list $(1)$ and negative index zero. In
dimension one, the list $(1,a)$ has one sign change exactly when $a<0$;
it has zero changes if $a=0$, as required for a zero eigenvalue.

In higher dimension write $a=M_{00}$. If $a\ne0$, eliminate the first
coordinate by a unit triangular congruence:
\[
 M\ \sim\ \diag(a,S),
\]
where $S$ is tridiagonal with first diagonal changed by $-b_0^2/a$,
and all remaining off-diagonal entries unchanged. It is again unreduced.
If $q_0=1,q_1,\ldots,q_{N-1}$ are its leading determinants, block
elimination on each leading block gives
\[
 (p_0,p_1,\ldots,p_N)=(1,a,aq_1,\ldots,aq_{N-1}).
\]
After deleting zeros, its variation is
$\boldsymbol1_{a<0}+V(q_0,\ldots,q_{N-1})$: the initial jump contributes
$\boldsymbol1_{a<0}$, and multiplying every later nonzero term by $a$
does not change successive sign comparisons. Congruence and induction
give the same recurrence for the negative index.

If $a=0$, unreducedness gives $b=b_0\ne0$. Use the first two coordinates
as one pivot block:
\[
 B=\begin{pmatrix}0&b\\b&c\end{pmatrix},\qquad
 B^{-1}=\begin{pmatrix}-c/b^2&1/b\\1/b&0\end{pmatrix}.
\]
The coupling to the remaining coordinates has zero first row, so
$E^\top B^{-1}E=0$. The block Schur complement is exactly the trailing
matrix $C=M_{2:,2:}$, without any update, and $M\sim\diag(B,C)$.
The block $B$ has one positive and one negative direction: its quadratic
form $2bst+ct^2$ becomes $2buv$ after $u=s+c t/(2b)$, $v=t$, and then
becomes a nonzero scalar times the difference of two squares.
Thus $\negind(M)=1+\negind(C)$.

If $q_0=1,q_1,\ldots,q_{N-2}$ are the leading determinants of $C$, the
same block elimination on every leading block gives
\[
 (p_0,p_1,\ldots,p_N)=(1,0,-b^2,-b^2q_1,\ldots,-b^2q_{N-2}).
\]
Deleting the zero contributes exactly one initial sign change, followed
by $V(q_0,\ldots,q_{N-2})$. Induction therefore gives the negative index
again. This includes the case of an empty trailing block. Finally apply
the conclusion to $M=T-xI$ and use the spectral interpretation of inertia.
\end{proof}

\paragraph{Why the unreduced assumption matters for deleting zeros.}
For $M=\diag(0,-1)$ the leading determinants are $(1,0,0)$, whose
zero-deleted variation is zero although $\negind(M)=1$. For a reduced
tridiagonal matrix, first split at each zero off-diagonal entry and add
the counts of its unreduced diagonal blocks. The finite block-diagonal
inertia identity justifies this procedure. This gives a zero-safe rule
without asserting an invalid zero-deletion rule for arbitrary matrices.

\section*{Handoff priorities}
\addcontentsline{toc}{section}{Handoff priorities}
All eight scaffold statements identified in this snapshot have complete
mathematical proofs above. The shortest bridges are the one-column norm
identity, the normalized rank-one projection identity, and the Givens
support invariant. Pivoted Cholesky should first establish its positive
prefix independently of rank. Hessenberg reduction should first expose
the padded-product identity. For the current Sturm scaffold, symmetric
LU followed by the existing inertia theorem is the shortest route;
the zero-safe block argument is a stronger optional extension.

The implicit-Q proof deliberately stays in the matrix layer. Its correct
intertwining orientation avoids an unjustified assumption on the second
matrix's subdiagonal. Every proof includes its zero or empty-dimensional
case where the statement permits one. No new statement has been marked
proved in the atlas, and no Lean source has been changed by this handoff.
\end{document}
